\documentclass[11pt,letterpaper]{article}
\usepackage[letterpaper,textwidth=437pt,textheight=634pt,hcentering,top=81pt,headsep=22pt,headheight=16pt]{geometry}
\usepackage[T1]{fontenc}
\usepackage{lmodern}
\usepackage{amsmath,amssymb,amsthm}
\usepackage[hidelinks]{hyperref}
\hypersetup{pdftitle={Full-cycle commutators in the Rubik's Cube group},
  pdfauthor={Ruimin Yan},pdfsubject={Finite permutation groups},
  pdfkeywords={Rubik's Cube group, commutator, permutation group, augmentation module, Lean}}

\theoremstyle{definition}
\newtheorem{theorem}{Theorem}[section]
\newtheorem{lemma}[theorem]{Lemma}
\newtheorem{proposition}[theorem]{Proposition}
\newtheorem{corollary}[theorem]{Corollary}
\theoremstyle{remark}
\newtheorem{example}[theorem]{Example}
\newcommand{\F}{\mathbb F}
\newcommand{\sgn}{\operatorname{sgn}}
\newcommand{\ord}{\operatorname{ord}}
\newcommand{\im}{\operatorname{im}}
\newcommand{\id}{\mathrm{id}}
\newcommand{\FC}{\mathcal F}
\newcommand{\cl}[2]{[#1,#2]}

\title{\fontsize{20}{24}\selectfont\bfseries Full-cycle commutators\\ in the Rubik's Cube group}
\author{Ruimin Yan}
\date{November 26, 2026}
\pagestyle{myheadings}
\markright{Full-cycle commutators in the Rubik's Cube group}

\begin{document}
\maketitle
\begin{abstract}
Every even position of the fixed-centre Rubik's Cube is a single commutator,
by a theorem of Timothy Sun. We show that both factors can be required to
induce an eight-cycle on the corners and a twelve-cycle on the edges.
Each factor then has exact order \(24\). We give two proofs. The first
deduces the factorization from a long-cycle theorem of Bader, B\"ohmler
and Wegener by lifting permutations through the corner and edge orientation
modules. For a fixed pair of edge cycles, the only obstruction to lifting
all orientations is that the cycles induce the same alternating
bipartition. The second proof, formalized in Lean~4, verifies finite
certificates for the position cycle types and treats pure edge flips
separately.
\end{abstract}

\section{The restricted commutator map}

Let \(G\) be the group of the ordinary \(3\times3\times3\) Rubik's Cube,
with the six centres fixed and unmarked. Write \(G^+\) for the subgroup
of positions inducing even permutations on both the corner slots and the
edge slots. We compose actions from right to left and put
\(\cl{A}{B}=ABA^{-1}B^{-1}\). The sign of the corner permutation is a
homomorphism on \(G\), so every commutator lies in \(G^+\). Sun proved the
converse: every element of \(G^+\) is a single commutator
\cite[Theorem~18]{sun}. We ask how far one may restrict the two factors
while retaining this image.

Let \(\FC\subseteq G\) consist of the positions whose corner permutation
is an \(8\)-cycle and whose edge permutation is a \(12\)-cycle. An element
of \(\FC\) moves every corner and every edge. The orientation increments
of its pieces are subject to the usual legality conditions.

\begin{theorem}\label{thm:main}
The map
\[
  \FC\times\FC\longrightarrow G^+,\qquad (A,B)\longmapsto\cl{A}{B},
\]
is surjective. Every member of \(\FC\) has order \(24\).
\end{theorem}

For example, a position which flips just two edges and leaves every
other piece fixed admits such a factorization. All position changes
cancel in the commutator, even though both factors move every piece.

Prescribing the cycle types of commutator factors is a classical question
for symmetric groups; see Vavpeti\v{c}~\cite{vavpetic}. Bader, B\"ohmler and
Wegener~\cite[Theorem~3.1 and Proposition~3.4]{bbw} give the permutation
input used in our first proof. Their signed-cycle theorem prescribes
negative sign product on each factor. A legal edge action on the cube
has positive sign product, since it flips an even number of edges.
We use their positional construction and their bipartition argument on
the even orientation space; the corner lift takes place in characteristic
three. Section~\ref{sec:cert} describes a second proof of the same
cube-group statement using explicit witnesses in dimensions eight and
twelve. The two proofs have different dependencies.

\section{Lifting a positional commutator}\label{sec:lift}

For \(m\in\{2,3\}\), let \(I_n=\{0,\ldots,n-1\}\) and set
\[
 V_{n,m}=\{v\in\F_m^{I_n}:\textstyle\sum_{i\in I_n}v_i=0\}.
\]
The pair \((p,v)\), with \(p\in S_n\) and \(v\in V_{n,m}\), sends an
oriented piece \((i,z)\) to \((p(i),z+v_i)\). Thus the increment is
indexed by the source slot. Write \(H_{n,m}\) for the group of these
actions. The standard cubie-coordinate theorem identifies \(G\) with
\begin{equation}\label{eq:cube}
 \big\{((p,v),(q,w))\in H_{8,3}\times H_{12,2}:
                       \sgn(p)=\sgn(q)\big\}.
\end{equation}
Here \(v\) records corner twists and \(w\) records edge flips. The two
zero-sum conditions and the equality of signs are also sufficient for
face-turn reachability \cite[\S3]{sun}.

For \(c\in S_n\) put \((\Delta_c x)_i=x_{c(i)}-x_i\). This map preserves
\(V_{n,m}\). Fix a target \(T=(t,u)\) and write \(A=(a,x)\), \(B=(b,y)\).
The equality \(\cl{A}{B}=T\) is equivalent to \(AB=TBA\). Applying both
sides to \((i,0)\) gives the two conditions
\begin{align}
 ab&=tba, \label{eq:position}\\
 \Delta_b x-\Delta_a y&=u\circ ba. \label{eq:orientation}
\end{align}
For example, the two increments before subtraction are
\(y_i+x_{b(i)}\) and \(x_i+y_{a(i)}+u_{b(a(i))}\). This also proves
sufficiency, since both actions add the initial orientation \(z\).
Once \(a,b\) have been chosen, the possible right-hand sides in
\eqref{eq:orientation} form the sum of the images of the two difference
maps restricted to \(V_{n,m}\).

\begin{lemma}\label{lem:corner}
If \(c\) is an \(n\)-cycle and \(m\nmid n\), then
\(\Delta_c:V_{n,m}\longrightarrow V_{n,m}\) is an isomorphism.
\end{lemma}
\begin{proof}
A vector fixed by an \(n\)-cycle is constant. A constant vector with value
\(k\) has coordinate sum \(nk\), which vanishes only for \(k=0\).
The claim follows by equality of dimensions.
\end{proof}

For even \(n\), let \(\mathcal P(c)=\{R_c,B_c\}\) be the two orbits of
\(c^2\) for an \(n\)-cycle \(c\). The two classes alternate around the
cycle. The following description of the binary difference map is the
balanced-bipartition lemma of \cite[Lemmas~5.4 and~5.6]{bbw}; we give the
argument in source-slot coordinates.

\begin{lemma}\label{lem:edge}
For even \(n\ge4\), in \(V=V_{n,2}\) one has
\[
 \im(\Delta_c|_V)=U_c,\qquad
 U_c=\{w\in V:\textstyle\sum_{i\in R_c}w_i=0\}.
\]
This is a hyperplane of \(V\). For \(n\)-cycles \(a,b\),
\(U_a+U_b=V\) if and only if \(\mathcal P(a)\ne\mathcal P(b)\).
\end{lemma}
\begin{proof}
Summing \(\Delta_c x\) over \(R_c\) gives the sum of \(x\) over all
slots in characteristic two. Thus the image is contained in \(U_c\).
A vector with ones at one red and one blue slot shows that \(U_c\) is a
proper hyperplane. The kernel on \(V\) is the one-dimensional space of
constant vectors, so the image and \(U_c\) both have dimension \(n-2\).
If two unordered bipartitions differ, two slots in the same class of
one lie in different classes of the other. Their two-point indicator
distinguishes the hyperplanes. Two distinct hyperplanes span \(V\);
identical bipartitions give the same proper hyperplane.
\end{proof}

\section{Proof of the factorization}\label{sec:proof}

We use the following consequence of \cite[Theorem~3.1 and
Proposition~3.4]{bbw}. For \(n\ge6\), each \(t\in A_n\) has a
representation \(t=\cl{a}{b}\) with \(a,b\) both \(n\)-cycles. If \(n\)
is even and \(t\ne\id\), one can moreover arrange
\(\mathcal P(a)\ne\mathcal P(b)\). For clarity, Proposition~3.4 there
gives a representative of each nonidentity even cycle type using the
standard cycle and a normalized parity-unstable cycle. A parity defect
in the latter places a label in the opposite alternating class, so the
two bipartitions differ. Their convention for the commutator is
\(a^{-1}b^{-1}ab\); replacing both cycles by their inverses converts it
to ours and preserves both bipartitions. Simultaneous conjugation then
gives the prescribed \(t\).

\begin{proposition}\label{prop:oneorbit}
For \((n,m)=(8,3)\) or \((12,2)\), every \((t,u)\in H_{n,m}\) with
\(t\in A_n\) is a commutator of two elements whose position permutations
are \(n\)-cycles.
\end{proposition}
\begin{proof}
For the corners, choose \(8\)-cycles \(a,b\) with \(t=\cl{a}{b}\).
Since \(3\nmid8\), Lemma~\ref{lem:corner} solves
\(\Delta_b x=u\circ ba\) in \(V_{8,3}\); take \(y=0\).
This includes \(t=\id\), for which one can take \(a=b\).

For the edges with \(t\ne\id\), choose \(12\)-cycles \(a,b\) with
distinct alternating bipartitions as above. Lemma~\ref{lem:edge} shows
that every balanced \(u\circ ba\) belongs to \(U_b+U_a\).
Consequently \eqref{eq:orientation} has a solution with \(x,y\in V_{12,2}\).

It remains to produce a pure flip \((\id,u)\). Pair the ones of \(u\)
among themselves, and the zeros among themselves. Both multiplicities
are even, so this gives six disjoint pairs, including when \(u=0\) or
\(u\) is all ones. Let \(R\) be the union of any three pairs. It has six
slots and \(\sum_{i\in R}u_i=0\). Choose a cycle \(c\) alternating between
\(R\) and its complement, and take \(a=b=c\). Its square preserves
\(R\), hence \(u\circ c^2\in U_c\). Lemma~\ref{lem:edge} now supplies
\(x\in V_{12,2}\) with \(\Delta_c x=u\circ c^2\); take \(y=0\).
This is the twelve-slot case of the support selection in
\cite[Lemma~6.6]{bbw}. In all three cases \eqref{eq:position}
and \eqref{eq:orientation} give the asserted commutator.
\end{proof}

\begin{proof}[Proof of Theorem~\ref{thm:main}]
Apply Proposition~\ref{prop:oneorbit} to the corner and edge components
of \(g\in G^+\). Write these components as \(\cl{A_c}{B_c}\) and
\(\cl{A_e}{B_e}\). The \(8\)-cycle in \(A_c\) and the \(12\)-cycle in
\(A_e\) both have negative sign. Their offsets are balanced, so
\(A=(A_c,A_e)\) belongs to \(G\) by \eqref{eq:cube}. The same holds for
\(B=(B_c,B_e)\). Commutators are taken componentwise in the product,
which gives \(\cl{A}{B}=g\).

If \((c,x)\in H_{n,m}\) has an \(n\)-cycle as its permutation, its
\(n\)-th power fixes each slot and adds \(\sum_i x_i=0\) at each slot.
Every power acting trivially on positions has exponent divisible by \(n\).
Thus the corner and edge actions of any member of \(\FC\) have exact
orders \(8\) and \(12\), respectively. The order in their direct product
is \(\operatorname{lcm}(8,12)=24\).
\end{proof}

\begin{corollary}[Sun's single-commutator theorem]
\(G^+\) is the derived subgroup of \(G\), and each of its elements is
a single commutator in \(G\).
\end{corollary}
\begin{proof}
The corner sign kills every commutator, hence its kernel contains the
derived subgroup. Theorem~\ref{thm:main} puts every element of that
kernel in the image of the commutator map.
\end{proof}

\begin{example}[The superflip]\label{ex:superflip}
Number the corner and edge slots cyclically, starting at zero, and let
\(c_n(i)=i+1\bmod n\). Give \(A\) and \(B\) the same two position cycles
\(c_8,c_{12}\). Give both zero corner offsets; on edges take \(x=0\)
for \(A\) and \(y_i=i\bmod2\) for \(B\). Exactly six entries of \(y\)
are nonzero. The left side of \eqref{eq:orientation} is the all-ones
edge vector, and the position commutator is trivial. Thus \(\cl{A}{B}\)
flips every edge in place and fixes the corners. This pair is explicit
in the accompanying formalization.
\end{example}

\section{An independent finite proof in Lean}\label{sec:cert}

The ancillary development proves Theorem~\ref{thm:main} in the group
\eqref{eq:cube}. It replaces the cited long-cycle construction by
finite witnesses and checks the orientation lifts directly. Its
declaration \texttt{CubeComm.every\_even\_cube} asserts membership of the
two witnesses in the legal group, their full cycles, their exact orders
and their commutator equation. The identification of \eqref{eq:cube}
with the face-turn group uses the classical reachability theorem.

For each position cycle type, choose a representative \(t\). A row
specifies full cycles \(a,b\) and a balanced solution of
\eqref{eq:orientation} for each target \(u=e_j-e_{n-1}\),
\(0\le j<n-1\), where \(e_j\) is a coordinate unit vector. These targets
form a basis of \(V_{n,m}\), including in characteristic two when \(m\)
divides \(n\). Indeed, any balanced vector \(u\) is
\(\sum_{j<n-1}u_j(e_j-e_{n-1})\). The linear combination of the stored
solutions with these coefficients therefore solves the equation for \(u\).
The number of basis witnesses checked is
\[
   \underbrace{12\cdot7}_{\text{corners}}+
   \underbrace{39\cdot11}_{\text{nonidentity edge types}}=513.
\]
For example, the corner identity type still has seven independent
twist targets. The edge identity type uses a separate argument.

Coverage follows from cycle type in the full symmetric group. The
catalogue includes the partitions \(\lambda\) of \(n\) for which
\(n-\ell(\lambda)\) is even, omitting \(1^{12}\) only on the edge side.
Lean checks these lists against the cycle types of the actual rows.
Conjugating by a relabelling \(d\) sends \((p,v)\) to
\((dpd^{-1},v\circ d^{-1})\), so both balance and \eqref{eq:orientation}
transport to any labelling of the same type. The relabelling of one
piece orbit need not itself be a physical move: it is a change of
coordinates, and the resulting paired factors satisfy \eqref{eq:cube}.

For pure edge flips, the formal proof uses three fixed six-element sets
\[
 P=\{0,1,2,3,4,5\},\qquad Q=\{0,1,2,6,7,8\},\qquad R=P\mathbin{\triangle}Q.
\]
At least one of their three subset sums vanishes in \(\F_2\), since the
third is the sum of the first two. For each set a fixed alternating
cycle and explicit offsets solve the equation whenever its sum vanishes.
This covers all \(2^{11}\) balanced flip vectors. Finally the two
permutation signs, the componentwise equality and the order calculation
give the result in the subgroup \(G\).

\paragraph{Reproducibility.}
The ancillary files include the Lean sources, the finite witnesses and
the dependency manifest. The optional generator searches for witnesses;
Lean checks each row, its cycle label and the exhaustiveness assertion.
With the pinned Lean~4.35.0 toolchain and mathlib revision in the
manifest, run the following commands from the ancillary directory:
\begin{verbatim}
bash cache.sh
lake build CubeComm
lake env lean -DwarningAsError=true AuditPaper.lean
lake env leanchecker CubeComm
\end{verbatim}
The audit checks the main theorem and its single-commutator corollary.
Their transitive axiom dependencies are propositional extensionality,
classical choice and quotient soundness.

\paragraph{Use of AI tools.}
OpenAI Codex assisted in the preparation of the Lean development,
the literature comparison, and the drafting of the mathematical exposition.

\begin{thebibliography}{9}
\bibitem{bbw} P.~Bader, B.~B\"ohmler, and P.~Wegener,
\emph{Commutators of signed \(n\)-cycles}, preprint, 2026.
\href{https://arxiv.org/abs/2609.24518v1}{arXiv:2609.24518v1 [math.GR]}.

\bibitem{sun} T.~Sun, \emph{Commutators in the Rubik's Cube group},
Amer. Math. Monthly \textbf{131} (2024), no.~1, 3--19.
\href{https://doi.org/10.1080/00029890.2023.2263158}{doi:10.1080/00029890.2023.2263158}.

\bibitem{vavpetic} A.~Vavpeti\v{c}, \emph{Commutators of cycles in permutation
groups}, Ars Math. Contemp. \textbf{10} (2016), no.~1, 67--77.
\href{https://doi.org/10.26493/1855-3974.430.eaf}{doi:10.26493/1855-3974.430.eaf}.
\end{thebibliography}
\end{document}
